%%%%%%%%%%%%%%%%%%%%%%%%%%%%%%%%%%%%%%%%%%%%%%%%%%%%%%%%%%%%%%%%%%%%%%
% Naslovnica, kazalo in številčenje za samostojni PDF poročila.
% main.tex ostane poglavje zbornika; to datoteko se vključi namesto
% uvodnih strani zbornika, tik pred \import{Skupina20/}{main}.
%%%%%%%%%%%%%%%%%%%%%%%%%%%%%%%%%%%%%%%%%%%%%%%%%%%%%%%%%%%%%%%%%%%%%%
\makeatletter

\hypersetup{
  pdftitle={Meja prepletenosti na procesorju IBM: koliko kubitov lahko zares prepletemo v stanje GHZ},
  pdfauthor={Dominik Krašovec, Erik Radovičevič}
}

% Številke razdelkov, enačb, slik in tabel brez številke poglavja.
\renewcommand\thesection{\@arabic\c@section}
\renewcommand\thesubsection{\thesection.\@arabic\c@subsection}
\renewcommand\theequation{\@arabic\c@equation}
\renewcommand\thefigure{\@arabic\c@figure}
\renewcommand\thetable{\@arabic\c@table}

% Naslov je na naslovnici, zato ga poglavje in glave strani ne ponovijo.
\def\title#1{}
\def\titlerunning#1{}
\def\author#1{}
\def\maketitle{}

% Brez praznih strani; številka je na vsaki strani, tudi na začetku literature.
\let\cleardoublepage\clearpage
\let\clearemptydoublepage\clearpage
\let\startnewpage\clearpage
\let\ps@bchap\ps@plain
\renewcommand\backmatter{\clearpage}

\makeatother

% Naslovnica
\thispagestyle{empty}
\begin{center}
\vspace*{2.5cm}
{\LARGE\bfseries Meja prepletenosti na procesorju IBM: koliko kubitov lahko zares prepletemo v stanje GHZ\par}
\vspace{1.2cm}
{\large Seminarska naloga za kvantno procesiranje pri predmetu Nekonvencionalne platforme in metode procesiranja\par}
\end{center}
\vspace{5cm}
\begin{flushright}
Avtorja: Dominik Krašovec, Erik Radovičevič
\end{flushright}
\vfill
\begin{center}
Ljubljana, November 2026
\end{center}
\clearpage

% Kazalo, brez številke strani
\tableofcontents
\thispagestyle{empty}
\clearpage

% Številčenje se začne takoj za kazalom.
\pagenumbering{arabic}
\pagestyle{plain}
